\documentclass[../../../include/open-logic-section]{subfiles}

\begin{document}

\olfileid{sth}{replacement}{absinf}
\olsection{الاستبدال و«اللانهاية المطلقة»}

سنترك الآن~\limofsize{} وراءنا، ونبحث عائلة مختلفة من المحاولات
(الداخلية) لتسويغ الاستبدال، تأخذ على محمل الجد فكرة تكوّن المجموعات
على مراحل.

عندما رسمنا العملية التكرارية أول مرة، قدمنا بعض المبادئ التي شرحت ما
يحدث في كل مرحلة. وكانت هي~\stageshier و\stagesord و\stagesacc. ثم
أضفنا لاحقًا مبادئ تخبرنا شيئًا عن عدد المراحل: فقد أخبرنا
\stagessucc{} أن عملية تكوين المجموعات لا تنتهي قط، وأخبرنا
\stagesinf{} أن العملية تمر بمرحلة لا نهائية.

ومن المعقول اقتراح أن هذين المبدأين الأخيرين ينتجان من مبدأ أعم، مثل:
\begin{enumerate}
	\item[] \stagesinex. عدد المراحل لا نهائي على نحو مطلق؛ والهرم
	مرتفع إلى أقصى ارتفاع ممكن.
\end{enumerate}
من الواضح أن هذا مبدأ غير صوري. لكن حتى إن لم يكن \emph{لازمًا}
مباشرة عن التصور التراكمي-التكراري للمجموعة، فإنه يبدو قطعًا
\emph{منسجمًا} معه. فهو، على أقل تقدير وبخلاف~\limofsize، يحتفظ بفكرة
تكوّن المجموعات مرحلة بعد مرحلة.

والأمل الآن هو توظيف~\stagesinex{} في تسويغ الاستبدال. فلنر كيف يمكن
فعل ذلك.

رأينا في~\olref[ordinals][idea]{sec} أنه يسهل بناء ترتيب جيد ينبغي،
بحسب المقصود، أن يكون متماثلًا مع~$\omega+\omega$.
وبعبارة أخرى، يسهل أن نتصور عملية تكرارية تجري مرحلة بعد مرحلة ويكون
نمط ترتيبها (بحسب المقصود)~$\omega+\omega$. ولذلك، إذا قبلنا
\stagesinex، فلا بد أن نقبل وجود المرحلة ذات الرتبة
$\omega+\omega$ على الأقل من الهرم، أي~$V_{\omega+\omega}$، إذ من
المؤكد أن الهرم \emph{يستطيع} الاستمرار إلى هذا الحد.

تتعمم هذه الفكرة هكذا: لكل ترتيب جيد، ينبغي أن تستمر عملية بناء الهرم
التكراري إلى حد لا يقل عن ذلك الترتيب الجيد. ويمكننا ضمان ذلك بمجرد
معاملة~\olref[ordinals][ordtype]{thmOrdinalRepresentation} بوصفها
\emph{مسلّمة}. وهذا سيخبرنا أن كل ترتيب جيد متماثل مع عدد ترتيبي لفون
نويمان. ولما كان كل عدد ترتيبي لفون نويمان مساويًا لرتبته الخاصة،
ستخبرنا~\olref[ordinals][ordtype]{thmOrdinalRepresentation} عندئذ
أنه كلما استطعنا وصف ترتيب جيد في نظرية المجموعات لدينا، وجب أن تتجاوز
عملية البناء التكراري للمجموعات ذلك الترتيب الجيد.

تبدو هذه الفكرة قطعًا نتيجة لـ\stagesinex. لكن إذا كانت غايتنا
استخراج الاستبدال منها، فإننا نواجه حاجزًا تقنيًا بسيطًا: الاستبدال
أقوى تمامًا من~\olref[ordinals][ordtype]{thmOrdinalRepresentation}.
(يبدي~\citet[\S13.2]{Potter2004} هذه الملاحظة؛ وسنثبتها في
\olref[sth][replacement][finiteaxiomatizability]{sec}.)

والخلاصة أنه إذا أردنا فهم~\stagesinex{} على نحو ينتج الاستبدال، فلا
يمكن أن يقول \emph{فحسب} إن الهرم يتجاوز كل ترتيب جيد؛ بل لا بد أن
يقرر ادعاءً أقوى. ولهذه الغاية، اقترح~\citet{Shoenfield:AST} تقوية
طبيعية جدًا للفكرة، هي أن الهرم ليس \emph{مسايرًا} لأي مجموعة.
\footnote{يبدو أن غودل اقترح فكرة مشابهة؛ انظر
  \citet[p.~223]{Potter2004}. ولمناقشة غودل و
  \citeauthor{Shoenfield:AST}، انظر~\citet[90--5]{Incurvati2020}.}
وبشيء من التفصيل: إذا كانت~$\tau$ تطبيقًا يرسل المجموعات إلى مراحل من
الهرم، فإن صورة أي مجموعة~$A$ تحت~$\tau$ لا تستنفد الهرم. وبعبارة
أخرى (مخططيًا):
\begin{enumerate}
	\item[] \stagescofin. إذا كانت~$A$ مجموعة وكانت~$\tau(x)$ مرحلة
	لكل~$x \in A$، فهناك مرحلة تأتي بعد كل~$\tau(x)$، حيث~$x \in A$.
\end{enumerate}
من الواضح أن~$\ZF$ تثبت صيغة مناسبة من~\stagescofin. وعلى العكس،
يمكننا أن نحاج غير صوريًا بأن~\stagescofin{} يسوّغ الاستبدال.
\footnote{سيكون إثبات الاستبدال باستعمال صياغة صورية ما لـ
  \stagescofin أصعب، لأن~$\Z$ وحدها ليست قوية بما يكفي لتعريف
  المراحل، ولذلك لا يتضح كيف يمكن صوغ~\stagescofin صوريًا. لكن أحد
  الخيارات هو العمل في امتداد ما لـ$\LT$، كما نوقش في
  \olref[sth][replacement][extrinsic]{sec}.}
فلتفترض~$(\forall x \in A)\lexists![y][\phi(x,y)]$. ولكل~$x \in A$،
لتكن~$\sigma(x)$ هي~$y$ التي تحقق~$\phi(x,y)$، ولتكن~$\tau(x)$ هي
المرحلة التي تتكون فيها~$\sigma(x)$ أول مرة. بموجب~\stagescofin، توجد
مرحلة~$V$ تحقق~$(\forall x \in A)\tau(x)\in V$. والآن، لما كان كل
$\tau(x) \in V$ وكانت~$\sigma(x) \subseteq \tau(x)$، ولما كانت~$V$
مقتدرة، كان~$\sigma(x) \in V$. ومن ثم نستطيع، باستعمال الفصل، الحصول
على
$\Setabs{y \in V}{(\exists x \in A)\sigma(x) = y} =
\Setabs{y}{(\exists x \in A)\phi(x,y)}$.

\begin{prob}
	صُغ~\stagescofin{} صوريًا داخل~$\ZF$.
\end{prob}

وهكذا تسوّغ~\stagescofin{} الاستبدال. ومن المعقول على الأقل أن تسوّغ
\stagesinex{} المبدأ~\stagescofin. فافترض أن~\stagescofin{} يفشل.
عندئذ يكون الهرم مسايرًا لمجموعة ما~$A$، أي لدينا تطبيق~$\tau$ يحقق
أنه لكل مرحلة~$S$ يوجد~$x \in A$ بحيث~$S \in \tau(x)$. وفي هذه
الحالة تكون لدينا بالفعل وسيلة للإحاطة بـ«اللانهاية المطلقة» المزعومة
للهرم: إذ \emph{يستنفدها} مدى~$\tau$ المطبق على~$A$. وهذا يضعف فكرة
أن الهرم «لا نهائي على نحو مطلق». وبأخذ المقابل بالنقيض: تستلزم
\stagesinex{} المبدأ~\stagescofin، الذي يسوّغ بدوره الاستبدال.

يمثل هذا محاولة واعدة حقًا لتقديم تسويغ داخلي للاستبدال. لكننا سنترك
لك الحكم فيما إذا كانت تنجح في نهاية المطاف.

\end{document}
